\documentclass[11pt, a4paper]{article}
\usepackage[utf8]{inputenc}
\usepackage[T1]{fontenc}
\usepackage[spanish]{babel}
\usepackage{amsmath, amssymb}
\usepackage{geometry}
\usepackage{enumitem}
\usepackage{titlesec}
\usepackage{parskip}
\usepackage{mdframed}
\usepackage{xcolor}

\geometry{margin=2.5cm}

\titleformat{\section}{\large\bfseries}{}{0em}{}
\titleformat{\subsection}{\normalsize\bfseries}{}{0em}{}
\titleformat{\subsubsection}{\normalsize\itshape}{}{0em}{}

\newmdenv[
    linecolor=gray!50,
    linewidth=0.5pt,
    backgroundcolor=gray!8,
    skipabove=6pt,
    skipbelow=6pt,
    innertopmargin=6pt,
    innerbottommargin=6pt,
]{respuesta}

\newcommand{\pregunta}[1]{\subsection{#1}}
\newcommand{\eje}[1]{\subsubsection{#1}}

\title{\textbf{Modelo de Parcial 1 -- Neurociencia y Comportamiento} \\
\large Inteligencia Artificial y Neurociencias \\ Universidad Torcuato Di Tella}
\author{}
\date{}

\begin{document}
\maketitle

\begin{mdframed}[linecolor=black, linewidth=0.7pt]
\textbf{Formato} (igual al simulacro): en el parcial caen \textbf{cuatro preguntas}, una por cada eje tem\'atico. Responder cada una en \textbf{entre uno y tres p\'arrafos, no m\'as}. Los ejes son: (1)~Modelado computacional del comportamiento humano, (2)~Aprendizaje, (3)~Emociones, (4)~Gen\'etica del comportamiento. Debajo de cada pregunta se da una \emph{respuesta posible} a modo de gu\'ia.
\end{mdframed}

\eje{1. Modelado computacional del comportamiento humano}
\pregunta{En el art\'iculo de Nature sobre profundidad de planificaci\'on en \textit{Four-in-a-Row}, el modelo predice tres tipos de datos observables en los humanos. \textquestiondown Cu\'ales son y qu\'e componente del modelo permite predecir cada uno?}

\begin{respuesta}
\textit{Respuesta posible.} El modelo predice tres cosas. Primero, las \textbf{elecciones} (jugadas) de las personas: validado con validaci\'on cruzada 5-fold en partidas humano vs.\ humano, predice las jugadas muy por encima del azar, gracias a la funci\'on de valor y al algoritmo de b\'usqueda que arman la variaci\'on principal. Segundo, los \textbf{tiempos de respuesta}: el tama\~no del \'arbol de decisi\'on construido por el algoritmo de b\'usqueda es predictor del tiempo de respuesta observado (\'arboles m\'as grandes, m\'as tiempo). Tercero, los \textbf{movimientos oculares}: la distribuci\'on de casillas que visita el algoritmo de b\'usqueda correlaciona con la distribuci\'on de la atenci\'on visual de los jugadores ($\rho \approx 0{,}54$). La conclusi\'on central es que los jugadores m\'as fuertes planifican m\'as profundamente y tienen menos lapsus atencionales, sin evidencia de que mejore la calidad de la heur\'istica (los pesos de las features).
\end{respuesta}

\eje{2. Aprendizaje}
\pregunta{\textquestiondown Cu\'ales son las tres condiciones que, seg\'un Darwin, deben cumplirse para que ocurra la selecci\'on natural? Explique adem\'as por qu\'e la teor\'ia de Lamarck no requiere una de ellas.}

\begin{respuesta}
\textit{Respuesta posible.} Para que ocurra selecci\'on natural deben darse tres condiciones: \textbf{variabilidad} (existen diferencias entre los rasgos de los individuos de una poblaci\'on), \textbf{heredabilidad} (esas variaciones se heredan de padres a hijos) y \textbf{adaptabilidad} (algunas de esas variaciones se adaptan mejor al ambiente, por lo que los individuos que las poseen sobreviven y se reproducen m\'as, dejando m\'as descendencia). Sin variaciones previas en la poblaci\'on, el proceso de selecci\'on no tiene sobre qu\'e operar.

En la teor\'ia de \textbf{Lamarck} los rasgos \emph{adquiridos} durante la vida del individuo se heredan a la descendencia; por eso no se necesita que existan variaciones previas en la poblaci\'on al comienzo del proceso: la variaci\'on se genera durante la vida y se transmite. En la de \textbf{Darwin}, en cambio, los rasgos adquiridos no se heredan, y las variaciones preexistentes (y heredables) son la materia prima imprescindible de la selecci\'on natural.
\end{respuesta}

\eje{3. Emociones}
\pregunta{Mencionamos cuatro afirmaciones sobre las emociones en las que la psicolog\'ia evolutiva s\'i est\'a de acuerdo. Enum\'erelas y explique brevemente la que vincula las emociones con la interocepci\'on.}

\begin{respuesta}
\textit{Respuesta posible.} Las cuatro afirmaciones son: (1)~las emociones son el resultado del \textbf{perfeccionamiento de dispositivos de supervivencia} a lo largo de la evoluci\'on de las especies; (2)~esos dispositivos usan las \textbf{mismas estructuras cerebrales y los mismos neurotransmisores en todos los mam\'iferos}, incluido el ser humano; (3)~esas estructuras est\'an vinculadas a \textbf{funciones b\'asicas de regulaci\'on del equilibrio interno}, necesario para la supervivencia; y (4)~las emociones est\'an \'intimamente arraigadas en fen\'omenos sensitivos viscerales corporales que llamamos \textbf{interocepci\'on}.

Sobre la cuarta: la interocepci\'on es la percepci\'on de los estados viscerales del organismo (pulsaciones, respiraci\'on, tensi\'on, estado del est\'omago, etc.). Las emociones no son ajenas al cuerpo, sino que est\'an enraizadas en esas se\~nales corporales. En la formulaci\'on extrema de William James, ``los cambios f\'isicos que experimentamos ante una situaci\'on \emph{son} las emociones'': la emoci\'on no antecede a la reacci\'on corporal, sino que en buena medida consiste en ella.
\end{respuesta}

\eje{4. Gen\'etica del comportamiento}
\pregunta{Explique el m\'etodo cl\'asico de comparaci\'on entre gemelos y mellizos para estimar el componente gen\'etico de un rasgo. \textquestiondown Qu\'e l\'ogica lo sustenta y cu\'al es su principal supuesto?}

\begin{respuesta}
\textit{Respuesta posible.} Los gemelos (monocig\'oticos) comparten pr\'acticamente el 100\,\% de sus alelos, mientras que los mellizos (dicig\'oticos) comparten en promedio la mitad. El m\'etodo compara la \textbf{correlaci\'on} de un rasgo entre pares de gemelos con la correlaci\'on entre pares de mellizos (criados juntos y del mismo sexo). Si la similitud entre gemelos es mayor que entre mellizos, se sospecha que el rasgo tiene un componente gen\'etico, porque la \'unica diferencia sistem\'atica entre ambos tipos de pares es cu\'antos alelos comparten. En f\'ormula:
\[
C_G = 2 \times [\mathrm{Corr}(\text{gemelos}) - \mathrm{Corr}(\text{mellizos})], \quad
C_{AC} = \mathrm{Corr}(\text{gemelos}) - C_G, \quad
C_{ANC} = 1 - C_G - C_{AC}.
\]

El \textbf{supuesto clave} es que los ambientes que generan similitudes son los mismos para gemelos que para mellizos (\emph{equal environments assumption}). Si se viola, hay sesgos: si los gemelos comparten ambientes m\'as parecidos que los mellizos (por ejemplo por tener a alguien id\'entico en el mundo), el m\'etodo \textbf{sobreestima} el componente gen\'etico; si los ambientes intrauterinos son m\'as similares entre mellizos que entre gemelos, lo \textbf{subestima}.
\end{respuesta}

\end{document}
